\paragraph{The rectangular condition and Kimura finiteness.}
There is a useful exact interpretation of condition~(ii).
\citet[Theorem~3.2, p.~286]{c_delpadronemazza2005}
prove that, under $C^+(X)$, vanishing of the minimal rectangular
Schur functor $S_{((b+1)^{a+1})}(h(X))$ is equivalent to
Kimura finite-dimensionality of $h(X)$, and to nilpotence of the
numerically trivial endomorphism ideal on every $h(X^m)$.
Their statement allows any chosen Weil cohomology.
Thus the finite criterion does not impose Kimura finiteness as
an additional hypothesis, but its rectangular condition implies
that property under $C^+(X)$.  Arbitrary Schur-finiteness cannot
replace the displayed rectangle.  Our proof uses the parity lift
only in a categorical trace and needs neither a Chow-idempotent
parity lift nor an unproved nilpotence theorem.
